\documentclass[12pt,a4paper]{article}

% Packages
\usepackage[utf8]{inputenc}
\usepackage[portuguese]{babel}
\usepackage{graphicx}
\usepackage{geometry}
\usepackage{fancyhdr}
\usepackage{listings}
\usepackage{xcolor}
\usepackage{hyperref}
\usepackage{tikz}
\usepackage{amsmath}
\usepackage{float}
\usepackage{caption}
\usepackage{enumitem}
\usepackage{booktabs}
\usepackage{longtable}

% Geometry
\geometry{left=2.5cm, right=2.5cm, top=3cm, bottom=3cm}

% Header/Footer
\pagestyle{fancy}
\fancyhf{}
\fancyhead[L]{Googol - Motor de Busca Distribuído}
\fancyhead[R]{\thepage}
\renewcommand{\headrulewidth}{0.4pt}

% Colors for code
\definecolor{codegreen}{rgb}{0,0.6,0}
\definecolor{codegray}{rgb}{0.5,0.5,0.5}
\definecolor{codepurple}{rgb}{0.58,0,0.82}
\definecolor{backcolour}{rgb}{0.95,0.95,0.92}

% Code listing style
\lstdefinestyle{mystyle}{
    backgroundcolor=\color{backcolour},   
    commentstyle=\color{codegreen},
    keywordstyle=\color{magenta},
    numberstyle=\tiny\color{codegray},
    stringstyle=\color{codepurple},
    basicstyle=\ttfamily\footnotesize,
    breakatwhitespace=false,         
    breaklines=true,                 
    captionpos=b,                    
    keepspaces=true,                 
    numbers=left,                    
    numbersep=5pt,                  
    showspaces=false,                
    showstringspaces=false,
    showtabs=false,                  
    tabsize=2
}
\lstset{style=mystyle}

% Hyperref setup
\hypersetup{
    colorlinks=true,
    linkcolor=blue,
    filecolor=magenta,      
    urlcolor=cyan,
    pdftitle={Googol - Relatório de Arquitetura},
    pdfpagemode=FullScreen,
}

\usetikzlibrary{shapes.geometric, arrows, positioning, fit, backgrounds}

\tikzstyle{component} = [rectangle, rounded corners, minimum width=3cm, minimum height=1cm, text centered, draw=black, fill=blue!20]
\tikzstyle{barrel} = [rectangle, rounded corners, minimum width=2.5cm, minimum height=1cm, text centered, draw=black, fill=green!20]
\tikzstyle{service} = [rectangle, rounded corners, minimum width=3cm, minimum height=1cm, text centered, draw=black, fill=orange!20]
\tikzstyle{arrow} = [thick,->,>=stealth]

\begin{document}

% Title Page
\begin{titlepage}
    \centering
    \vspace*{2cm}
    
    {\Huge\bfseries Googol\par}
    \vspace{0.5cm}
    {\LARGE Motor de Busca Distribuído\par}
    \vspace{1cm}
    {\Large\bfseries Relatório de Arquitetura do Sistema\par}
    
    \vspace{2cm}
    
    {\Large Sistemas Distribuídos\par}
    \vspace{0.5cm}
    {\large Universidade de Coimbra\par}
    {\large Departamento de Engenharia Informática\par}
    
    \vfill
    
    {\large\bfseries Autores:\par}
    \vspace{0.5cm}
    \begin{tabular}{ll}
        Henrique Diz & 2023213681 \\
        Rodrigo Manão & 2023207589 \\
        João Francisco & 2023228417 \\
    \end{tabular}
    
    \vspace{1cm}
    
    {\large \today\par}
\end{titlepage}

% Table of Contents
\tableofcontents
\newpage

% Abstract
\begin{abstract}
Este relatório descreve a arquitetura completa do Googol, um motor de busca distribuído desenvolvido como projeto da disciplina de Sistemas Distribuídos. O sistema utiliza tecnologias como Java RMI para comunicação entre componentes distribuídos, Spring Boot para a camada de API REST, e React/Next.js para o frontend. A arquitetura implementa conceitos avançados como balanceamento de carga, tolerância a falhas, replicação de dados, deteção dinâmica de stopwords, e análise contextual com inteligência artificial.
\end{abstract}

\newpage

% ========================================
% SECTION 1: INTRODUÇÃO
% ========================================
\section{Introdução}

\subsection{Visão Geral}
O Googol é um motor de busca distribuído que implementa uma arquitetura escalável e tolerante a falhas. O sistema é composto por múltiplos componentes que comunicam através de Java RMI (Remote Method Invocation), coordenados por um Gateway central que gere o balanceamento de carga e a agregação de resultados.

\subsection{Objetivos do Sistema}
\begin{itemize}
    \item Indexação distribuída de páginas web
    \item Pesquisa eficiente com suporte a múltiplas palavras
    \item Balanceamento de carga entre servidores de índice (Barrels)
    \item Tolerância a falhas com replicação de dados
    \item Interface web moderna e responsiva
    \item Análise contextual com IA
    \item Estatísticas em tempo real
\end{itemize}

\subsection{Tecnologias Utilizadas}
\begin{itemize}
    \item \textbf{Backend:} Java 21, Spring Boot 3.2.0, Java RMI, JSoup
    \item \textbf{Frontend:} React 19, Next.js 16, TypeScript, Framer Motion, GSAP, WebGL (OGL)
    \item \textbf{Comunicação:} RMI Registry, REST API (HTTP/JSON)
    \item \textbf{IA:} OpenRouter API (integração com modelos de linguagem)
\end{itemize}

\newpage

% ========================================
% SECTION 2: ARQUITETURA GLOBAL
% ========================================
\section{Arquitetura Global do Sistema}

\subsection{Diagrama de Arquitetura}

\begin{figure}[H]
\centering
\begin{tikzpicture}[node distance=2cm]

% Frontend
\node (frontend) [component] {Frontend (React/Next.js)};

% WebApp
\node (webapp) [service, below of=frontend, yshift=-0.5cm] {WebApp (Spring Boot)};

% Gateway
\node (gateway) [component, below of=webapp, yshift=-0.5cm] {Gateway (Port 8183)};

% Barrels
\node (barrel1) [barrel, below left of=gateway, xshift=-3cm, yshift=-1cm] {Barrel 1\\Port 8182};
\node (barrel2) [barrel, below of=gateway, yshift=-1cm] {Barrel 2\\Port 8184};
\node (barrel3) [barrel, below right of=gateway, xshift=3cm, yshift=-1cm] {Barrel 3\\Port 8185};

% Queue
\node (queue) [component, below of=barrel2, yshift=-1.5cm] {URL Queue\\Port 8181};

% Downloader
\node (downloader) [service, below of=queue, yshift=-0.5cm] {Downloader(s)};

% Arrows
\draw [arrow] (frontend) -- node[anchor=west] {HTTP/JSON} (webapp);
\draw [arrow] (webapp) -- node[anchor=west] {RMI} (gateway);
\draw [arrow] (gateway) -- (barrel1);
\draw [arrow] (gateway) -- (barrel2);
\draw [arrow] (gateway) -- (barrel3);
\draw [arrow] (barrel1) -- (queue);
\draw [arrow] (barrel2) -- (queue);
\draw [arrow] (barrel3) -- (queue);
\draw [arrow] (queue) -- (downloader);
\draw [arrow, bend left=60] (downloader) to node[anchor=west] {Index} (barrel1);
\draw [arrow] (downloader) -- (barrel2);
\draw [arrow, bend right=60] (downloader) to (barrel3);

\end{tikzpicture}
\caption{Arquitetura global do sistema Googol}
\label{fig:architecture}
\end{figure}

\subsection{Componentes Principais}

\begin{enumerate}
    \item \textbf{Frontend:} Interface web em React/Next.js
    \item \textbf{WebApp:} API REST em Spring Boot (porta 8080)
    \item \textbf{Gateway:} Coordenador central RMI (porta 8183)
    \item \textbf{Barrels:} Servidores de índice invertido (portas 8182, 8184, 8185)
    \item \textbf{URL Queue:} Fila de URLs a indexar (porta 8181)
    \item \textbf{Downloader(s):} Serviços de download e processamento de páginas
\end{enumerate}

\subsection{Fluxo de Dados}

\subsubsection{Fluxo de Indexação}
\begin{enumerate}
    \item Utilizador submete URL através do frontend
    \item WebApp recebe pedido e envia para Gateway via RMI
    \item Gateway verifica se URL já foi indexado
    \item Se não, Gateway adiciona URL à URL Queue
    \item Downloader retira URL da fila (operação bloqueante)
    \item Downloader faz download da página com JSoup
    \item Downloader extrai palavras e links
    \item Downloader envia palavras para TODOS os Barrels ativos
    \item Cada Barrel armazena no índice invertido e retransmite para outros Barrels
    \item Downloader adiciona links descobertos à fila (baixa prioridade)
\end{enumerate}

\subsubsection{Fluxo de Pesquisa}
\begin{enumerate}
    \item Utilizador introduz query no frontend
    \item Frontend envia pedido GET para WebApp
    \item WebApp chama Gateway.searchWords() via RMI
    \item Gateway divide query em palavras
    \item Para cada palavra, Gateway chama Barrel.searchWord() (round-robin)
    \item Gateway faz interseção dos resultados (AND lógico)
    \item Gateway obtém contagens de inbound links de todos os Barrels
    \item Gateway ordena resultados por número de referências (PageRank-like)
    \item Gateway aplica paginação
    \item WebApp enriquece resultados com PageInfo (título, descrição)
    \item WebApp retorna JSON para frontend
    \item Frontend exibe resultados com AnimatedList
\end{enumerate}

\newpage

% ========================================
% SECTION 3: BACKEND - RMI
% ========================================
\section{Backend - Camada RMI}

\subsection{Gateway}

\subsubsection{Responsabilidades}
O Gateway é o componente central do sistema, responsável por:
\begin{itemize}
    \item \textbf{Balanceamento de carga:} Round-robin entre Barrels ativos
    \item \textbf{Tolerância a falhas:} Retry automático em caso de falha de Barrel
    \item \textbf{Coordenação de pesquisa:} Agregação e interseção de resultados
    \item \textbf{Gestão de Barrels:} Registo, ativação/desativação, sincronização
    \item \textbf{Integração com fila:} Verificação de URLs já indexados
    \item \textbf{Estatísticas:} Coleta de métricas de pesquisa e tempos de resposta
    \item \textbf{Notificações em tempo real:} Callbacks para atualizações de estado
\end{itemize}

\subsubsection{Interface RMI}
\begin{lstlisting}[language=Java, caption={GatewayInterface - Métodos principais}]
public interface GatewayInterface extends Remote {
    // Pesquisa
    List<String> searchWordGateway(String word) throws RemoteException;
    List<String[]> searchWords(List<String> words, int page, int pageSize) 
        throws RemoteException;
    int getTotalResults(List<String> words) throws RemoteException;
    
    // Gestao de URLs
    String addURL(String url, boolean indexAnyway) throws RemoteException;
    List<String[]> getUrlsForIndexedUrl(String url) throws RemoteException;
    
    // Gestao de Barrels
    void registerBarrel(String host, int port, String name) 
        throws RemoteException;
    void unActiveBarrel(String name, int port) throws RemoteException;
    List<String> getActiveBarrels() throws RemoteException;
    List<String> getRegisteredBarrels() throws RemoteException;
    
    // Estatisticas e Callbacks
    Map<String, Object> getBarrelStatistics() throws RemoteException;
    void registerStatsCallback(StatsCallback callback) 
        throws RemoteException;
    void registerBarrelStateCallback(BarrelStateCallback callback) 
        throws RemoteException;
}
\end{lstlisting}

\subsubsection{Algoritmo de Balanceamento de Carga}
\begin{lstlisting}[language=Java, caption={Round-robin load balancing}]
private BarrelInterface getNextBarrel() {
    if (activeBarrels.isEmpty()) return null;
    BarrelInterface barrel = activeBarrels.get(currentBarrelIndex);
    currentBarrelIndex = (currentBarrelIndex + 1) % activeBarrels.size();
    return barrel;
}
\end{lstlisting}

\subsubsection{Algoritmo de Pesquisa Multi-palavra}
\begin{lstlisting}[language=Java, caption={Interseção de resultados para múltiplas palavras}]
// Obter resultados para cada palavra
List<HashSet<String>> resultsPerWord = new ArrayList<>();
for (String word : words) {
    resultsPerWord.add(new HashSet<>(searchWordGateway(word)));
}

// Intersecao (AND logico)
HashSet<String> intersection = new HashSet<>(resultsPerWord.get(0));
for (int i = 1; i < resultsPerWord.size(); i++) {
    intersection.retainAll(resultsPerWord.get(i));
}

// Ranking por inbound links
List<String[]> ranked = new ArrayList<>();
for (String url : intersection) {
    int refCount = inboundLinkCounts.getOrDefault(url, 0);
    ranked.add(new String[]{url, String.valueOf(refCount)});
}

// Ordenacao decrescente
ranked.sort((a, b) -> Integer.compare(
    Integer.parseInt(b[1]), 
    Integer.parseInt(a[1])
));
\end{lstlisting}

\subsection{Barrels (Servidores de Índice)}

\subsubsection{Estruturas de Dados}
\begin{lstlisting}[language=Java, caption={Estruturas principais dos Barrels}]
// Indice invertido: palavra -> conjunto de URLs
private ConcurrentHashMap<String, HashSet<String>> indexedItems;

// Grafo de URLs: URL -> conjunto de URLs linkados
private ConcurrentHashMap<String, HashSet<String>> urlsIndexed;

// Prevencao de duplicacao
private Set<String> receivedMessages;

// Gestor de stopwords
private BarrelStopWords stopWordsManager;

// Backup da fila
private BlockingDeque<String> queueBackup;
private Set<String> queueSeenUrls;
\end{lstlisting}

\subsubsection{Replicação de Índice}
Quando um Barrel recebe uma entrada de índice, retransmite para outros Barrels:

\begin{lstlisting}[language=Java, caption={Retransmissão para replicação}]
public void addToIndex(String word, String url, boolean shouldRetransmit) 
    throws RemoteException {
    // Adicionar ao indice local
    indexedItems.computeIfAbsent(word, k -> new HashSet<>()).add(url);
    
    // Retransmitir para outros Barrels
    if (shouldRetransmit) {
        List<String> activeBarrels = gateway.getActiveBarrels();
        ExecutorService executor = Executors.newCachedThreadPool();
        
        for (String barrelKey : activeBarrels) {
            if (!barrelKey.equals(myKey)) {
                executor.submit(() -> {
                    BarrelInterface barrel = getBarrelByKey(barrelKey);
                    barrel.addToIndex(word, url, false); // Nao retransmitir novamente
                });
            }
        }
    }
}
\end{lstlisting}

\subsection{Deteção Dinâmica de Stopwords}

\subsubsection{Algoritmo de Deteção}
O sistema utiliza análise estatística para identificar stopwords dinamicamente:

\begin{enumerate}
    \item \textbf{Contagem de frequências:} Para cada página, contar ocorrências de palavras
    \item \textbf{Deteção de outliers:} Usar método IQR (Interquartile Range)
    \begin{align*}
        Q1 &= \text{Percentil 25} \\
        Q3 &= \text{Percentil 75} \\
        IQR &= Q3 - Q1 \\
        \text{Threshold} &= Q3 + (2.0 \times IQR)
    \end{align*}
    \item \textbf{Análise cross-document:} Palavra torna-se stopword se aparecer como outlier em $\geq 70\%$ dos documentos
    \item \textbf{Sincronização:} Quando stopwords mudam, retransmitir para todos os Barrels
\end{enumerate}

\begin{lstlisting}[language=Java, caption={Configuração de stopwords}]
private static final double STOPWORD_THRESHOLD = 0.70;  // 70% documentos
private static final double OUTLIER_K = 2.0;             // Multiplicador IQR
private static final int MIN_WORD_FREQUENCY = 5;         // Frequencia minima
\end{lstlisting}

\subsection{URL Queue}

\subsubsection{Gestão de Prioridades}
A fila implementa duas prioridades:
\begin{itemize}
    \item \textbf{Alta prioridade:} URLs adicionados manualmente pelo utilizador (frente da fila)
    \item \textbf{Baixa prioridade:} URLs descobertos pelo Downloader (fim da fila)
\end{itemize}

\begin{lstlisting}[language=Java, caption={Implementação de fila com prioridade}]
private BlockingDeque<String> urlsToIndex;  // Deque para prioridade
private Set<String> seenUrls;                // Prevencao de duplicados

public void putNew(String url, boolean priority) throws RemoteException {
    if (!seenUrls.add(url)) return;  // Ja visto
    if (priority) 
        urlsToIndex.putFirst(url);   // Alta prioridade
    else 
        urlsToIndex.putLast(url);    // Baixa prioridade
}
\end{lstlisting}

\subsection{Downloader}

\subsubsection{Pipeline de Processamento}
\begin{enumerate}
    \item Obter URL da fila (operação bloqueante)
    \item Download da página com JSoup
    \item Tokenização do conteúdo textual
    \item Para cada palavra:
    \begin{itemize}
        \item Converter para minúsculas
        \item Remover caracteres não alfanuméricos
        \item Filtrar palavras com $< 3$ caracteres
        \item Contar frequência
        \item Enviar para todos os Barrels ativos via \texttt{addToIndex()}
    \end{itemize}
    \item Extrair todos os links \texttt{<a href>}
    \item Adicionar links à fila (baixa prioridade)
    \item Enviar backlinks para Barrels via \texttt{addUrlsForIndexedUrl()}
    \item Enviar frequências para análise de stopwords via \texttt{addWordCounts()}
\end{enumerate}

\subsubsection{Retry Logic}
O Downloader implementa retry com backoff:
\begin{itemize}
    \item 3 tentativas por URL
    \item Timeout de 10 segundos por tentativa
    \item Em caso de falha persistente, URL é descartado
\end{itemize}

\subsection{Persistência e Recuperação}

\subsubsection{Backup de Estado}
Os Barrels guardam estado em disco (\texttt{barrelX\_progress.dat}):
\begin{itemize}
    \item Índice invertido (\texttt{indexedItems})
    \item Grafo de URLs (\texttt{urlsIndexed})
    \item Backup da fila (\texttt{urlsToIndex}, \texttt{seenUrls})
\end{itemize}

\subsubsection{Recuperação após Falha}
\begin{enumerate}
    \item Ao iniciar, Barrel verifica existência de ficheiro de progresso
    \item Se existe: carregar de disco
    \item Se não: copiar índice de outros Barrels ativos
    \item Registar no Gateway
    \item Sincronizar dados em falta
\end{enumerate}

\newpage

% ========================================
% SECTION 4: BACKEND - WebApp
% ========================================
\section{Backend - WebApp (Spring Boot)}

\subsection{Arquitetura em Camadas}

\begin{figure}[H]
\centering
\begin{tikzpicture}[node distance=1.5cm]
\node (controllers) [component] {Controllers};
\node (service) [service, below of=controllers] {Service Layer};
\node (rmi) [component, below of=service] {RMI Layer};

\draw [arrow] (controllers) -- node[anchor=west] {DTOs} (service);
\draw [arrow] (service) -- node[anchor=west] {RMI Calls} (rmi);
\end{tikzpicture}
\caption{Arquitetura em camadas do WebApp}
\end{figure}

\subsection{REST API Endpoints}

\begin{longtable}{|p{3cm}|p{3.5cm}|p{7cm}|}
\hline
\textbf{Endpoint} & \textbf{Método} & \textbf{Descrição} \\
\hline
\endfirsthead
\hline
\textbf{Endpoint} & \textbf{Método} & \textbf{Descrição} \\
\hline
\endhead
\hline
\endfoot

\texttt{/api/search} & GET & Pesquisa simples com query params \\
\texttt{/api/search} & POST & Pesquisa avançada com JSON body \\
\texttt{/api/connections} & GET & Obter backlinks de um URL \\
\texttt{/api/barrels/active} & GET & Listar Barrels ativos \\
\texttt{/api/barrels/registered} & GET & Listar Barrels registados \\
\texttt{/api/barrels/add-url} & POST & Adicionar URL ao índice \\
\texttt{/api/statistics} & GET & Obter estatísticas do sistema \\
\texttt{/api/health} & GET & Health check \\
\texttt{/api/info} & GET & Informação do sistema \\
\texttt{/api/context-analysis} & POST & Análise contextual com IA \\
\hline
\caption{Endpoints da REST API}
\label{tab:endpoints}
\end{longtable}

\subsection{DTOs (Data Transfer Objects)}

\subsubsection{SearchResponseDTO}
\begin{lstlisting}[language=Java, caption={Estrutura de resposta de pesquisa}]
public class SearchResponseDTO {
    private String query;
    private List<SearchResultDTO> results;
    private int currentPage;
    private int pageSize;
    private int totalResults;
    private int totalPages;
    private boolean hasResults;
}
\end{lstlisting}

\subsubsection{SearchResultDTO}
\begin{lstlisting}[language=Java, caption={Estrutura de resultado individual}]
public class SearchResultDTO {
    private String url;
    private String title;
    private String description;
    private int references;  // Numero de inbound links
}
\end{lstlisting}

\subsection{Configuração UTF-8}

Para garantir codificação correta de caracteres especiais (acentos, cedilhas):

\begin{lstlisting}[caption={application.properties}]
# UTF-8 Encoding
spring.http.encoding.charset=UTF-8
spring.http.encoding.enabled=true
spring.http.encoding.force=true
spring.http.encoding.force-request=true
spring.http.encoding.force-response=true

# JSON Encoding
spring.jackson.charset=UTF-8
\end{lstlisting}

\begin{lstlisting}[language=Java, caption={EncodingConfig.java}]
@Configuration
public class EncodingConfig implements WebMvcConfigurer {
    @Override
    public void configureMessageConverters(
        List<HttpMessageConverter<?>> converters) {
        StringHttpMessageConverter stringConverter = 
            new StringHttpMessageConverter(StandardCharsets.UTF_8);
        stringConverter.setWriteAcceptCharset(false);
        converters.add(0, stringConverter);
    }
}
\end{lstlisting}

\subsection{Integração RMI via Spring}

\begin{lstlisting}[language=Java, caption={RMIConfig.java - Bean de Gateway}]
@Configuration
public class RMIConfig {
    @Bean
    public GatewayInterface gatewayService() throws Exception {
        ConfigReader config = new ConfigReader("gateway");
        Registry registry = LocateRegistry.getRegistry(
            config.getHost(), 
            config.getPort()
        );
        return (GatewayInterface) registry.lookup(config.getName());
    }
}
\end{lstlisting}

\subsection{Análise Contextual com IA}

O sistema integra com OpenRouter API para análise contextual:

\begin{lstlisting}[language=Java, caption={ContextAnalysisController.java}]
@PostMapping("/context-analysis")
public ResponseEntity<Map<String, Object>> getContextAnalysis(
    @RequestBody Map<String, String> request) {
    
    String query = request.get("query");
    String citations = request.get("citations");
    
    String prompt = String.format(
        "Analisa a seguinte pesquisa: '%s'. " +
        "Contexto dos resultados: %s", 
        query, citations
    );
    
    // Chamada ao OpenRouter API
    Map<String, Object> response = openRouterService.analyze(prompt);
    return ResponseEntity.ok(response);
}
\end{lstlisting}

\newpage

% ========================================
% SECTION 5: FRONTEND
% ========================================
\section{Frontend - React/Next.js}

\subsection{Stack Tecnológica}

\begin{itemize}
    \item \textbf{Framework:} Next.js 16 (React 19) - App Router
    \item \textbf{Linguagem:} TypeScript (strict mode)
    \item \textbf{Animações:} Framer Motion, GSAP
    \item \textbf{WebGL:} OGL (shader 3D para orbe animado)
    \item \textbf{UI:} HeroUI, componentes custom
\end{itemize}

\subsection{Estrutura de Rotas}

\begin{longtable}{|p{4cm}|p{9cm}|}
\hline
\textbf{Rota} & \textbf{Descrição} \\
\hline
\endfirsthead
\hline
\textbf{Rota} & \textbf{Descrição} \\
\hline
\endhead
\hline
\endfoot

\texttt{/} & Página inicial com interface de pesquisa \\
\texttt{/results} & Resultados de pesquisa paginados \\
\texttt{/indexar} & Formulário para adicionar URLs \\
\texttt{/ligacoes} & Formulário para pesquisar backlinks \\
\texttt{/ligacoes/results} & Resultados de backlinks \\
\texttt{/statistics} & Dashboard de estatísticas em tempo real \\
\texttt{/autores} & Página da equipa \\
\hline
\caption{Rotas do frontend}
\end{longtable}

\subsection{Componentes Principais}

\subsubsection{AnimatedList}
Lista animada de resultados com funcionalidades:
\begin{itemize}
    \item Navegação por teclado (setas, Tab, Enter)
    \item Expansão/colapso de detalhes
    \item Auto-scroll para item selecionado
    \item Gradientes de fade no topo e base
    \item Animações com Framer Motion
\end{itemize}

\subsubsection{Orb (WebGL)}
Orbe 3D animado com shaders GLSL:
\begin{itemize}
    \item Simplex noise para distorção
    \item Rotação em hover
    \item Mudança de cor por hue (0-360°)
    \item Efeito wave baseado na posição do rato
    \item 60 FPS com requestAnimationFrame
\end{itemize}

\begin{lstlisting}[language=JavaScript, caption={Shader GLSL do Orb (fragmento)}]
vec3 color1 = hsv2rgb(vec3(hue / 360.0, 0.8, 1.0));
vec3 color2 = hsv2rgb(vec3((hue + 60.0) / 360.0, 0.9, 0.8));

float noise = snoise(vPosition * noiseScale + uTime * 0.2);
vec3 finalColor = mix(color1, color2, noise * 0.5 + 0.5);
\end{lstlisting}

\subsubsection{Cursor Customizado}
Cursor com trail de 30 círculos:
\begin{itemize}
    \item Física de spring (30\% damping)
    \item Cores adaptativas por rota
    \item Z-index alto (99999999)
    \item Esconde cursor padrão
\end{itemize}

\subsubsection{StaggeredMenu}
Menu lateral com animações GSAP:
\begin{itemize}
    \item 3 camadas de background com entrada staggered
    \item Itens com rotação + fade-in
    \item Efeito de texto alternado (Menu $\leftrightarrow$ Close)
    \item Rotação de ícone (0° → 225°)
\end{itemize}

\subsubsection{ProfileCard}
Cards 3D com efeito holográfico:
\begin{itemize}
    \item Tilt baseado na posição do rato
    \item Perspetiva 3D com CSS transforms
    \item Suporte a orientação do dispositivo (mobile)
    \item Links para redes sociais
    \item Efeito de shine holográfico
\end{itemize}

\subsection{Sistema de Cores por Rota}

\begin{longtable}{|p{3.5cm}|p{3cm}|p{3cm}|p{2cm}|}
\hline
\textbf{Rota} & \textbf{Cor Primária} & \textbf{Cor Secundária} & \textbf{Hue} \\
\hline
\endfirsthead
\hline
\textbf{Rota} & \textbf{Cor Primária} & \textbf{Cor Secundária} & \textbf{Hue} \\
\hline
\endhead
\hline
\endfoot

Home/Results & \texttt{\#9c43ff} (Roxo) & \texttt{\#4cb8e9} (Azul) & 0° \\
Indexar & \texttt{\#ff6b9d} (Rosa) & \texttt{\#ff8c42} (Laranja) & 270° \\
Ligações & \texttt{\#00ff88} (Verde) & \texttt{\#88ff00} (Amarelo) & 117° \\
Estatísticas & \texttt{\#ff3333} (Vermelho) & \texttt{\#cc0000} (Verm. Escuro) & 0° \\
\hline
\caption{Sistema de cores adaptativo}
\end{longtable}

\subsection{Funcionalidades Avançadas}

\subsubsection{Análise Contextual com IA}
\begin{itemize}
    \item Botão info flutuante ao lado da barra de pesquisa
    \item Hover para mostrar análise
    \item Cache de 30 minutos por query
    \item Envia query + top 5 descrições para backend
    \item Exibe insights contextuais
\end{itemize}

\subsubsection{Estatísticas em Tempo Real}
\begin{itemize}
    \item Polling a cada 2 segundos
    \item Top 10 pesquisas mais frequentes
    \item Tempo médio de resposta por Barrel
    \item Estado de Barrels (ativos/registados)
    \item Dashboard responsivo
\end{itemize}

\subsubsection{Paginação Inteligente}
\begin{itemize}
    \item Mostra primeiras 3, últimas 3, e atual ±2 páginas
    \item Ellipsis para páginas omitidas
    \item Sincronização com URL (\texttt{?page=N})
    \item Scroll suave ao mudar de página
\end{itemize}

\subsection{Otimizações de Performance}

\begin{enumerate}
    \item \textbf{Code splitting:} Automático via Next.js
    \item \textbf{Animações:} \texttt{requestAnimationFrame} para 60 FPS
    \item \textbf{Cache de API:} Análise contextual em cache (30 min)
    \item \textbf{Debouncing:} Polling de estatísticas a cada 2s
    \item \textbf{WebGL:} Contexto único, cleanup em unmount
    \item \textbf{Memory management:} Cleanup em todos os \texttt{useEffect}
\end{enumerate}

\subsection{Acessibilidade}

\begin{itemize}
    \item HTML semântico (\texttt{<main>}, \texttt{<nav>}, \texttt{<section>})
    \item Labels ARIA em todos os botões de ícone
    \item Navegação completa por teclado
    \item Estados de focus visíveis
    \item Alto contraste (branco em fundo escuro)
    \item Suporte a screen readers
\end{itemize}

\newpage

% ========================================
% SECTION 6: COMUNICAÇÃO
% ========================================
\section{Comunicação Entre Componentes}

\subsection{RMI (Remote Method Invocation)}

\subsubsection{Configuração de Registry}
\begin{lstlisting}[caption={Configuração RMI (Config.properties)}]
gateway.host=127.0.0.1
gateway.port=8183
gateway.name=gateway

queue.host=127.0.0.1
queue.port=8181
queue.name=urlqueue

barrel1.host=127.0.0.1
barrel1.port=8182
barrel1.name=barrel1

barrel2.host=127.0.0.1
barrel2.port=8184
barrel2.name=barrel2

barrel3.host=127.0.0.1
barrel3.port=8185
barrel3.name=barrel3
\end{lstlisting}

\subsubsection{Padrão de Callbacks}
O sistema implementa callbacks para notificações em tempo real:

\begin{lstlisting}[language=Java, caption={StatsCallback interface}]
public interface StatsCallback extends Remote {
    void onStatsUpdate(Map<String, Object> stats) throws RemoteException;
}
\end{lstlisting}

Uso:
\begin{enumerate}
    \item Cliente regista callback no Gateway
    \item Gateway armazena callback em \texttt{CopyOnWriteArrayList}
    \item Em mudança de estado, Gateway chama \texttt{callback.onStatsUpdate()}
    \item Cliente recebe e exibe estatísticas atualizadas
    \item Callbacks mortos removidos automaticamente em erro
\end{enumerate}

\subsection{REST API (HTTP/JSON)}

\subsubsection{Exemplo de Request/Response}

\textbf{Request:}
\begin{lstlisting}[language=bash]
GET /api/search?q=distributed%20systems&page=0&pageSize=10
\end{lstlisting}

\textbf{Response:}
\begin{lstlisting}[language=JSON]
{
  "query": "distributed systems",
  "results": [
    {
      "url": "https://example.com",
      "title": "Introduction to Distributed Systems",
      "description": "Learn about distributed computing...",
      "references": 42
    }
  ],
  "currentPage": 0,
  "pageSize": 10,
  "totalResults": 5,
  "totalPages": 1,
  "hasResults": true
}
\end{lstlisting}

\subsection{CORS (Cross-Origin Resource Sharing)}

\begin{lstlisting}[language=Java, caption={Configuração CORS}]
@Configuration
public class CorsConfig {
    @Bean
    public CorsFilter corsFilter() {
        CorsConfiguration config = new CorsConfiguration();
        config.setAllowCredentials(true);
        config.addAllowedOrigin("http://localhost:3000");
        config.addAllowedOrigin("http://localhost:3001");
        config.addAllowedHeader("*");
        config.addAllowedMethod("*");
        
        UrlBasedCorsConfigurationSource source = 
            new UrlBasedCorsConfigurationSource();
        source.registerCorsConfiguration("/api/**", config);
        
        return new CorsFilter(source);
    }
}
\end{lstlisting}

\newpage

% ========================================
% SECTION 7: ALGORITMOS CHAVE
% ========================================
\section{Algoritmos e Padrões Implementados}

\subsection{Algoritmos de Pesquisa}

\subsubsection{Pesquisa Multi-palavra com Interseção}
\begin{enumerate}
    \item Normalizar query (lowercase, split por espaços)
    \item Para cada palavra $w_i$:
    \begin{itemize}
        \item $R_i = \text{searchWord}(w_i)$ via Barrel (round-robin)
    \end{itemize}
    \item Calcular interseção: $R = R_1 \cap R_2 \cap \ldots \cap R_n$
    \item Obter inbound link counts de todos os Barrels
    \item Ordenar por \texttt{references} (descendente)
    \item Aplicar paginação: $\text{results}[\text{page} \times \text{pageSize} : (\text{page}+1) \times \text{pageSize}]$
\end{enumerate}

\subsubsection{Ranking por PageRank-like}
\begin{lstlisting}[language=Java, caption={Algoritmo de ranking}]
// Agregar contagens de inbound links
Map<String, Integer> inbound = new HashMap<>();
for (BarrelInterface barrel : activeBarrels) {
    Map<String, Integer> m = barrel.getInboundLinkCounts();
    for (Map.Entry<String, Integer> e : m.entrySet()) {
        inbound.merge(e.getKey(), e.getValue(), Integer::sum);
    }
}

// Ordenar por contagem (decrescente)
ranked.sort((a, b) -> Integer.compare(
    Integer.parseInt(b[1]),  // Inbound count de b
    Integer.parseInt(a[1])   // Inbound count de a
));
\end{lstlisting}

\subsection{Padrões de Design Implementados}

\begin{enumerate}
    \item \textbf{Distributed Hash Table (DHT):} Cada Barrel mantém parte do índice
    \item \textbf{Replication:} Entradas de índice replicadas em todos os Barrels
    \item \textbf{Load Balancing:} Round-robin entre Barrels
    \item \textbf{Failover:} Retry automático em falha de Barrel
    \item \textbf{Observer Pattern:} Callbacks para atualizações em tempo real
    \item \textbf{Service Locator:} RMI Registry para descoberta de serviços
    \item \textbf{DTO Pattern:} Separação entre objetos RMI e REST
    \item \textbf{Layered Architecture:} WebApp → Gateway → Barrels → Queue
    \item \textbf{Eventual Consistency:} Sincronização de índice via retransmissão
    \item \textbf{Persistence:} Estado guardado em disco para recuperação
\end{enumerate}

\newpage

% ========================================
% SECTION 8: TOLERÂNCIA A FALHAS
% ========================================
\section{Tolerância a Falhas e Resiliência}

\subsection{Mecanismos de Failover}

\subsubsection{Failover em Pesquisa}
\begin{lstlisting}[language=Java, caption={Retry automático}]
private List<String> searchWithFailover(String word) {
    for (int attempt = 0; attempt < activeBarrels.size(); attempt++) {
        try {
            BarrelInterface barrel = getNextBarrel();
            return barrel.searchWord(word);
        } catch (RemoteException e) {
            // Tentar proximo Barrel
            continue;
        }
    }
    // Todos falharam
    return Collections.emptyList();
}
\end{lstlisting}

\subsubsection{Recuperação de Barrel}
Quando um Barrel volta online:
\begin{enumerate}
    \item Barrel carrega estado de disco (se disponível)
    \item Regista-se no Gateway
    \item Gateway chama \texttt{syncIndex()} para sincronizar dados em falta
    \item Barrel recebe índice completo de outros Barrels
    \item Barrel marca-se como ativo
\end{enumerate}

\subsection{Backup e Recuperação}

\subsubsection{Backup de Fila}
Quando a URL Queue vai fazer shutdown:
\begin{lstlisting}[language=Java, caption={Backup para Barrels}]
public void shutdown() {
    List<BarrelInterface> barrels = getActiveBarrels();
    for (BarrelInterface barrel : barrels) {
        barrel.backupQueueState(urlsToIndex, seenUrls);
    }
}
\end{lstlisting}

\subsubsection{Restore de Fila}
Ao iniciar, URL Queue tenta restaurar de Barrels:
\begin{lstlisting}[language=Java, caption={Restore de backup}]
public void startup() {
    List<BarrelInterface> barrels = getActiveBarrels();
    for (BarrelInterface barrel : barrels) {
        Map<String, Object> backup = barrel.restoreQueueState();
        if (backup != null) {
            urlsToIndex = (BlockingDeque) backup.get("queue");
            seenUrls = (Set) backup.get("seen");
            return;
        }
    }
    // Sem backup, iniciar vazio
}
\end{lstlisting}

\subsection{Replicação de Dados}

\subsubsection{Consistência Eventual}
\begin{itemize}
    \item Downloader envia palavras para TODOS os Barrels
    \item Cada Barrel retransmite para outros Barrels
    \item Uso de \texttt{receivedMessages} previne duplicação
    \item Convergência garantida após todas as retransmissões
\end{itemize}

\subsubsection{Sincronização de Stopwords}
\begin{lstlisting}[language=Java, caption={Sincronização de stopwords}]
private void retransmitStopwords() {
    List<String> activeBarrels = gateway.getActiveBarrels();
    List<String> stopwords = new ArrayList<>(this.stopwords);
    
    for (String barrelKey : activeBarrels) {
        if (!barrelKey.equals(myKey)) {
            BarrelInterface barrel = getBarrelByKey(barrelKey);
            barrel.syncStopwords(stopwords);
        }
    }
}
\end{lstlisting}

\newpage

% ========================================
% SECTION 9: ESTATÍSTICAS E MONITORIZAÇÃO
% ========================================
\section{Estatísticas e Monitorização}

\subsection{Métricas Coletadas}

\begin{enumerate}
    \item \textbf{Estatísticas de Pesquisa:}
    \begin{itemize}
        \item \texttt{Map<String, Integer>} - Palavra → Contagem
        \item Top 10 pesquisas mais frequentes
    \end{itemize}
    
    \item \textbf{Tempo de Resposta:}
    \begin{itemize}
        \item \texttt{Map<String, Long>} - Barrel → Tempo total (ns)
        \item \texttt{Map<String, Long>} - Barrel → Contagem de requests
        \item Média = Tempo total / Contagem
    \end{itemize}
    
    \item \textbf{Estado de Barrels:}
    \begin{itemize}
        \item Lista de Barrels ativos
        \item Lista de Barrels registados
        \item Tamanho de índice por Barrel
    \end{itemize}
\end{enumerate}

\subsection{Dashboard de Estatísticas}

\begin{figure}[H]
\centering
\begin{tikzpicture}[node distance=1.5cm]
\node (top10) [component, minimum width=5cm, minimum height=3cm] {Top 10 Searches};
\node (tempo) [component, right of=top10, xshift=5cm, minimum width=5cm, minimum height=3cm] {Avg Response Time};
\node (barrels) [component, below of=top10, xshift=2.5cm, minimum width=10cm, minimum height=4cm] {Barrels Status};
\end{tikzpicture}
\caption{Layout do dashboard de estatísticas}
\end{figure}

\subsection{Atualização em Tempo Real}

\subsubsection{Polling do Frontend}
\begin{lstlisting}[language=JavaScript, caption={Polling a cada 2 segundos}]
useEffect(() => {
  const fetchStats = async () => {
    const stats = await getStatistics();
    const active = await getActiveBarrels();
    const registered = await getRegisteredBarrels();
    setData(stats);
    // ...
  };
  
  fetchStats(); // Inicial
  const interval = setInterval(fetchStats, 2000); // 2s
  return () => clearInterval(interval);
}, []);
\end{lstlisting}

\subsubsection{Callbacks RMI (Backend)}
\begin{lstlisting}[language=Java, caption={Notificação assíncrona}]
private void notifyStatsUpdateAsync() {
    for (StatsCallback cb : statsCallbacks) {
        callbackExecutor.submit(() -> {
            try {
                cb.onStatsUpdate(stats.getStatistics());
            } catch (Exception e) {
                statsCallbacks.remove(cb); // Auto-cleanup
            }
        });
    }
}
\end{lstlisting}

\newpage

% ========================================
% SECTION 10: CONCLUSÃO
% ========================================
\section{Conclusão}

\subsection{Objetivos Alcançados}

O projeto Googol implementou com sucesso:
\begin{itemize}
    \item \textbf{Arquitetura distribuída} com múltiplos componentes RMI
    \item \textbf{Balanceamento de carga} round-robin entre Barrels
    \item \textbf{Tolerância a falhas} com retry automático e replicação
    \item \textbf{Índice invertido} distribuído e replicado
    \item \textbf{Deteção dinâmica de stopwords} com análise estatística
    \item \textbf{Ranking por PageRank} baseado em inbound links
    \item \textbf{API REST} moderna com Spring Boot
    \item \textbf{Frontend responsivo} com animações avançadas e WebGL
    \item \textbf{Análise contextual com IA} via OpenRouter
    \item \textbf{Estatísticas em tempo real} com dashboard interativo
\end{itemize}

\subsection{Destaques Técnicos}

\subsubsection{Backend}
\begin{itemize}
    \item Comunicação RMI robusta com failover
    \item Replicação assíncrona de índice
    \item Persistência com serialização Java
    \item Callbacks para notificações em tempo real
    \item Integração Spring Boot + RMI
\end{itemize}

\subsubsection{Frontend}
\begin{itemize}
    \item WebGL shaders customizados (orbe 3D)
    \item Animações GSAP com timelines complexas
    \item Cursor customizado com física de spring
    \item Cards holográficos 3D
    \item Sistema de cores adaptativo por rota
    \item Paginação inteligente com ellipsis
\end{itemize}

\subsection{Pontos Fortes}

\begin{enumerate}
    \item \textbf{Escalabilidade:} Arquitetura permite adicionar Barrels dinamicamente
    \item \textbf{Resiliência:} Sistema continua operacional com falha de componentes
    \item \textbf{Performance:} Índice distribuído permite pesquisas paralelas
    \item \textbf{UX excepcional:} Interface moderna com animações fluidas
    \item \textbf{Inovação:} Deteção automática de stopwords e análise com IA
\end{enumerate}

\subsection{Possíveis Melhorias Futuras}

\begin{enumerate}
    \item \textbf{Sharding:} Distribuir índice por hash de palavra (não replicação total)
    \item \textbf{Caching:} Adicionar camada de cache Redis para queries frequentes
    \item \textbf{Ranking:} Implementar TF-IDF e algoritmo completo de PageRank
    \item \textbf{Crawling:} Downloader com rate limiting e robots.txt
    \item \textbf{Frontend:} Server-side rendering para SEO
    \item \textbf{Segurança:} Autenticação e rate limiting na API
    \item \textbf{Monitorização:} Prometheus + Grafana para métricas
\end{enumerate}

\subsection{Considerações Finais}

O Googol demonstra uma implementação completa e funcional de um motor de busca distribuído, combinando conceitos avançados de sistemas distribuídos (RMI, replicação, balanceamento) com tecnologias web modernas (Spring Boot, React, WebGL). O sistema é robusto, escalável e fornece uma excelente experiência ao utilizador.

A arquitetura modular permite fácil extensão e manutenção, enquanto os mecanismos de failover e replicação garantem alta disponibilidade. A interface moderna com animações WebGL e análise contextual com IA elevam o projeto além de um simples motor de busca académico.

\newpage

% ========================================
% REFERÊNCIAS
% ========================================
\section*{Referências}
\addcontentsline{toc}{section}{Referências}

\begin{enumerate}
    \item Oracle. \textit{Java RMI Tutorial}. \url{https://docs.oracle.com/javase/tutorial/rmi/}
    \item Spring Framework. \textit{Spring Boot Reference Documentation}. \url{https://spring.io/projects/spring-boot}
    \item Next.js. \textit{Next.js Documentation}. \url{https://nextjs.org/docs}
    \item React. \textit{React Documentation}. \url{https://react.dev/}
    \item GSAP. \textit{GreenSock Animation Platform}. \url{https://greensock.com/gsap/}
    \item OGL. \textit{OGL - Minimal WebGL Library}. \url{https://github.com/oframe/ogl}
    \item JSoup. \textit{Java HTML Parser}. \url{https://jsoup.org/}
    \item OpenRouter. \textit{AI Model Routing API}. \url{https://openrouter.ai/}
\end{enumerate}

\newpage

% ========================================
% APÊNDICES
% ========================================
\appendix

\section{Configuração do Ambiente}

\subsection{Requisitos}
\begin{itemize}
    \item Java 21
    \item Maven 3.8+
    \item Node.js 18+
    \item npm 9+
\end{itemize}

\subsection{Instalação}

\subsubsection{Backend}
\begin{lstlisting}[language=bash]
cd backend
mvn clean install
\end{lstlisting}

\subsubsection{Frontend}
\begin{lstlisting}[language=bash]
cd frontend
npm install
\end{lstlisting}

\subsection{Execução}

\subsubsection{Componentes RMI}
\begin{lstlisting}[language=bash]
# Gateway
cd backend
mvn exec:java -Dexec.mainClass="gateway.Gateway"

# Barrel 1
mvn exec:java -Dexec.mainClass="barrel.IndexStorageBarrel" \
  -Dexec.args="barrel1"

# Barrel 2
mvn exec:java -Dexec.mainClass="barrel.IndexStorageBarrel" \
  -Dexec.args="barrel2"

# URL Queue
mvn exec:java -Dexec.mainClass="queue.URLQueue"

# Downloader
mvn exec:java -Dexec.mainClass="downloader.Downloader"
\end{lstlisting}

\subsubsection{WebApp (Spring Boot)}
\begin{lstlisting}[language=bash]
cd backend
mvn spring-boot:run
\end{lstlisting}

\subsubsection{Frontend}
\begin{lstlisting}[language=bash]
cd frontend
npm run dev
\end{lstlisting}

\subsection{Scripts PowerShell (Windows Terminal)}
\begin{lstlisting}[language=bash]
# Backend com multiplos terminais
.\scripts\runBackendWT.ps1

# Frontend
.\scripts\runFrontendWT.ps1
\end{lstlisting}

\section{Estrutura de Ficheiros}

\subsection{Backend}
\begin{lstlisting}
backend/
├── src/main/java/
│   ├── barrel/
│   │   ├── IndexStorageBarrel.java
│   │   ├── BarrelInterface.java
│   │   ├── BarrelStopWords.java
│   │   └── BarrelProgress.java
│   ├── gateway/
│   │   ├── Gateway.java
│   │   └── GatewayInterface.java
│   ├── queue/
│   │   ├── URLQueue.java
│   │   └── URLQueueInterface.java
│   ├── downloader/
│   │   └── Downloader.java
│   ├── statistics/
│   │   ├── Statistics.java
│   │   ├── StatsCallback.java
│   │   └── BarrelStateCallback.java
│   ├── webapp/
│   │   ├── WebApplication.java
│   │   ├── controller/
│   │   ├── service/
│   │   ├── dto/
│   │   └── config/
│   └── common/
├── src/main/resources/
│   ├── Config.properties
│   └── application.properties
└── pom.xml
\end{lstlisting}

\subsection{Frontend}
\begin{lstlisting}
frontend/
├── src/
│   ├── app/
│   │   ├── layout.tsx
│   │   ├── page.tsx
│   │   ├── results/page.tsx
│   │   ├── indexar/page.tsx
│   │   ├── ligacoes/
│   │   ├── statistics/page.tsx
│   │   └── autores/page.tsx
│   ├── components/
│   │   ├── AnimatedList/
│   │   ├── Orb/
│   │   ├── Cursor/
│   │   ├── StaggeredMenu/
│   │   └── ProfileCard/
│   └── services/
│       └── api.ts
├── public/
├── package.json
└── tsconfig.json
\end{lstlisting}

\end{document}
